% Recommended, but optional, packages for figures and better typesetting:
\usepackage{microtype}
\usepackage{graphicx}
\usepackage{subfigure}
\newcommand{\theHalgorithm}{\arabic{algorithm}}
\usepackage{caption}
\usepackage{xspace}
\usepackage{pifont}
\usepackage{booktabs} % for professional tables
% hyperref makes hyperlinks in the resulting PDF.
% If your build breaks (sometimes temporarily if a hyperlink spans a page)
% please comment out the following usepackage line and replace
\usepackage{hyperref}
\usepackage{xcolor}
\usepackage{listings}
\usepackage{times}
\usepackage{latexsym}

\usepackage{algorithm}
\usepackage{algorithmic}

% For proper rendering and hyphenation of words containing Latin characters (including in bib files)
\usepackage[T1]{fontenc}
% For Vietnamese characters
% \usepackage[T5]{fontenc}
% See https://www.latex-project.org/help/documentation/encguide.pdf for other character sets

% This assumes your files are encoded as UTF8
\usepackage[utf8]{inputenc}
% Attempt to make hyperref and algorithmic work together better:



% Use the following line for the initial blind version submitted for review:
\usepackage[final]{acl}

% If accepted, instead use the following line for the camera-ready submission:

% For theorems and such
\usepackage{amsmath}
\usepackage{amssymb}
\usepackage{mathtools}
\usepackage{amsthm}
\usepackage{multirow}
\usepackage{graphicx}
\usepackage{lscape}
% if you use cleveref..
\usepackage[capitalize,noabbrev]{cleveref}
\newcommand{\xin}[1]{\textcolor{orange}{Xin: #1}}  % 
%%%%%%%%%%%%%%%%%%%%%%%%%%%%%%%%
% THEOREMS
%%%%%%%%%%%%%%%%%%%%%%%%%%%%%%%%
\theoremstyle{plain}
\newtheorem{theorem}{Theorem}[section]
\newtheorem{proposition}[theorem]{Proposition}
\newtheorem{lemma}[theorem]{Lemma}
\newtheorem{corollary}[theorem]{Corollary}
\theoremstyle{definition}
\newtheorem{definition}[theorem]{Definition}
\newtheorem{assumption}[theorem]{Assumption}
\theoremstyle{remark}
\newtheorem{remark}[theorem]{Remark}
\newcommand{\sysname}{\texttt{AgentSpec}\xspace}
% \newcommand{\sysnamew}{\texttt{SVDLLM-W}\xspace}
% \newcommand{\sysnamewl}{\texttt{SVDLLM-WL}\xspace}
\newcommand{\remaining}[1]{\textcolor{red}{Mark: #1}}

\definecolor{mygreen}{HTML}{009901}
\newcommand{\spup}[1]{\textcolor{mygreen}{#1}}
\newcommand{\spdown}[1]{\textcolor{red}{#1}}
\newcommand{\spone}[1]{\textcolor{gray}{#1}}
% Todonotes is useful during development; simply uncomment the next line
%    and comment out the line below the next line to turn off comments
%\usepackage[disable,textsize=tiny]{todonotes}
\usepackage[textsize=tiny]{todonotes}
\lstset{
  basicstyle=\ttfamily\small,
  columns=fullflexible,
  frame=single,
  breaklines=true,
  showstringspaces=false,
  keywordstyle=\color{blue},
  commentstyle=\spup,
  stringstyle=\color{blue},
}
\newcommand{\eeTable}{{
\begin{table*}[t]
\centering
\caption{Goodput (tokens/sec) of \sysname and baselines on four different agentic workloads with three different models. The best performance is marked in bold. The speedup is measured by comparing the goodput with the normal autoregressive decoding. The relative performance gain compared to the best-performing baseline is marked in green inside bracket.}
\label{tab:e2e_comparison}
\resizebox{0.95\linewidth}{!}{
\begin{tabular}{c|c|cc|cc|cc|cc}
\toprule
\textsc{Model} & \textsc{Method}
& \multicolumn{2}{c|}{Deep Research Agent}
& \multicolumn{2}{c|}{Code Generation Agent}
& \multicolumn{2}{c|}{GAIA}
& \multicolumn{2}{c}{SWE-Bench}\\
\cmidrule(lr){3-4}\cmidrule(lr){5-6}\cmidrule(lr){7-8}\cmidrule(lr){9-10}
& & Tokens/s & Speedup & Tokens/s & Speedup & Tokens/s & Speedup & Tokens/s & Speedup\\
\midrule

% ---------------- Qwen-3-8B ----------------
\multirow{5}{*}{\rotatebox[origin=c]{90}{\textbf{Qwen-3-8B}}} & Normal
& 617.41 & \spone{$\times$ 1.00}
& 636.09 & \spone{$\times$ 1.00}
& 501.27 & \spone{$\times$ 1.00}
& 652.44 & \spone{$\times$ 1.00} \\ \cmidrule{2-10}
         & EAGLE-3
& 605.14 & \spdown{$\times$ 0.98}
& 539.01 & \spdown{$\times$ 0.85}
& 468.22 & \spdown{$\times$ 0.93}
& 523.12 & \spdown{$\times$ 0.80} \\
         & NGram
& 509.62 & \spdown{$\times$ 0.83}
& 422.58 & \spdown{$\times$ 0.66}
& 308.27 & \spdown{$\times$ 0.61}
& 408.28 & \spdown{$\times$ 0.63} \\
         & SuffixDecoding
& 598.74 & \spdown{$\times$ 0.97}
& 579.21 & \spdown{$\times$ 0.91}
& 498.16 & \spdown{$\times$ 0.99}
& 556.64 & \spdown{$\times$ 0.85} \\ \cmidrule{2-10}
         & \sysname
& \textbf{661.43} \spup{($\uparrow 9\%$)}  & \spup{\textbf{$\times$ 1.07}}
& \textbf{828.10} \spup{($\uparrow 43\%$)} & \spup{\textbf{$\times$ 1.30}}
& \textbf{537.78} \spup{($\uparrow 8\%$)}  & \spup{\textbf{$\times$ 1.07}}
& \textbf{854.69} \spup{($\uparrow 42\%$)} & \spup{\textbf{$\times$ 1.31}} \\
\midrule

% ---------------- GPT-OSS-20B ----------------
\multirow{5}{*}{\rotatebox[origin=c]{90}{\textbf{GPT-OSS-20B}}} & Normal
& 201.43 & \spone{$\times$ 1.00}
& 289.26 & \spone{$\times$ 1.00}
& 247.67 & \spone{$\times$ 1.00}
& 351.22 & \spone{$\times$ 1.00} \\ \cmidrule{2-10}
            & EAGLE-3
& 125.71 & \spdown{$\times$ 0.62}
& 161.34 & \spdown{$\times$ 0.56}
& 112.16 & \spdown{$\times$ 0.45}
& 198.81 & \spdown{$\times$ 0.57} \\
            & NGram
& 158.91 & \spdown{$\times$ 0.79}
& 175.19 & \spdown{$\times$ 0.61}
& 109.22 & \spdown{$\times$ 0.44}
& 253.36 & \spdown{$\times$ 0.72} \\
            & SuffixDecoding
& 196.54 & \spdown{$\times$ 0.98}
& 295.85 & \spup{$\times$ 1.02}
& 187.79 & \spdown{$\times$ 0.76}
& 319.73 & \spdown{$\times$ 0.91} \\ \cmidrule{2-10}
            & \sysname
& \textbf{297.54} \spup{($\uparrow 51\%$)} & \spup{\textbf{$\times$ 1.48}}
& \textbf{584.43} \spup{($\uparrow 104\%$)} & \spup{\textbf{$\times$ 2.02}}
& \textbf{295.76} \spup{($\uparrow 57\%$)} & \spup{\textbf{$\times$ 1.19}}
& \textbf{595.20} \spup{($\uparrow 86\%$)} & \spup{\textbf{$\times$ 1.69}} \\
\midrule

% ---------------- DSK-8B ----------------
\multirow{5}{*}{\rotatebox[origin=c]{90}{\textbf{DSK-Distill-8B}}} & Normal
& 987.66  & \spone{$\times$ 1.00}
& 1230.28 & \spone{$\times$ 1.00}
& 823.39  & \spone{$\times$ 1.00}
& 1432.27 & \spone{$\times$ 1.00} \\ \cmidrule{2-10}
       & EAGLE-3
& 965.43  & \spdown{$\times$ 0.98}
& 1119.29 & \spdown{$\times$ 0.91}
& 864.43  & \spup{$\times$ 1.05}
& 1098.82 & \spdown{$\times$ 0.77} \\
       & NGram
& 752.66  & \spdown{$\times$ 0.76}
& 797.18  & \spdown{$\times$ 0.65}
& 725.29  & \spdown{$\times$ 0.88}
& 1014.47 & \spdown{$\times$ 0.71} \\
       & SuffixDecoding
& 995.14  & \spup{$\times$ 1.01}
& 1041.29 & \spdown{$\times$ 0.85}
& 820.09  & \spone{$\times$ 1.00}
& 1498.69 & \spup{$\times$ 1.05} \\ \cmidrule{2-10}
       & \sysname
& \textbf{1294.65} \spup{($\uparrow 30\%$)} & \spup{\textbf{$\times$ 1.31}}
& \textbf{1671.34} \spup{($\uparrow 49\%$)} & \spup{\textbf{$\times$ 1.36}}
& \textbf{924.47}  \spup{($\uparrow 7\%$)}  & \spup{\textbf{$\times$ 1.12}}
& \textbf{2281.87} \spup{($\uparrow 52\%$)} & \spup{\textbf{$\times$ 1.59}}\\
\bottomrule
\end{tabular}
}
\end{table*}
}}




% ================= Table: MiMo-7B USACO (Thpt + Spd) =================
\newcommand{\mimoTable}{{
\begin{table*}[ht]
\centering
\caption{Goodput (tokens/sec) of \sysname and baselines, including MTP, on MiMo-7B. The best performance is marked in bold. The speedup is measured by comparing the goodput with the normal autoregressive decoding. The relative performance gain compared to the best-performing baseline is marked in green inside bracket.}
\label{tab:mimo_usaco}
\resizebox{1\linewidth}{!}{
\begin{tabular}{c| c| cc cc cc cc cc}
\toprule
\textsc{Model} & \textsc{Method}
& \multicolumn{2}{c}{Code Gen.(Bronze)}
& \multicolumn{2}{c}{Code Gen.(Silver)}
& \multicolumn{2}{c}{Code Gen.(Gold)}
& \multicolumn{2}{c}{Code Gen.(Platinum)}
& \multicolumn{2}{c}{Code Gen.(All)} \\
\cmidrule(lr){3-4}\cmidrule(lr){5-6}\cmidrule(lr){7-8}\cmidrule(lr){9-10}\cmidrule(lr){11-12}
& & Tokens/s & Speedup & Tokens/s & Speedup & Tokens/s & Speedup & Tokens/s & Speedup & Tokens/s & Speedup \\
\midrule

\multirow{5}{*}{\rotatebox[origin=c]{90}{\textbf{MiMo-7B}}} & Normal
& 2768.67 & \spone{$\times$ 1.00}
& 3093.50 & \spone{$\times$ 1.00}
& 2532.30 & \spone{$\times$ 1.00}
& 1141.64 & \spone{$\times$ 1.00}
& 2572.16 & \spone{$\times$ 1.00} \\ \cmidrule{2-12}
        & NGram
& 1322.92 & \spdown{$\times$ 0.48}
& 1399.76 & \spdown{$\times$ 0.45}
& 1383.92 & \spdown{$\times$ 0.55}
& 990.87   & \spdown{$\times$ 0.87}
& 1306.92 & \spdown{$\times$ 0.51} \\
        & SuffixDecoding
& 1609.41 & \spdown{$\times$ 0.58}
& 1624.98 & \spdown{$\times$ 0.53}
& 1599.63 & \spdown{$\times$ 0.63}
& 1235.12 & \spup{$\times$ 1.08}
& 1653.26 & \spdown{$\times$ 0.64} \\
        & MTP
& 1924.46 & \spdown{$\times$ 0.70}
& 2001.28 & \spdown{$\times$ 0.65}
& 2034.30 & \spdown{$\times$ 0.80}
& 1349.10 & \spup{$\times$ 1.18}
& 1875.18 & \spdown{$\times$ 0.73} \\ \cmidrule{2-12}
        & \sysname
& \textbf{3324.71} \spup{($\uparrow 73\%$)} & \spup{$\times$ \textbf{1.20}}
& \textbf{3751.44} \spup{($\uparrow 87\%$)} & \spup{$\times$ \textbf{1.21}}
& \textbf{2925.40} \spup{($\uparrow 44\%$)} & \spup{$\times$ \textbf{1.16}}
& \textbf{1527.14} \spup{($\uparrow 13\%$)} & \spup{$\times$ \textbf{1.34}}
& \textbf{3318.99} \spup{($\uparrow 77\%$)} & \spup{$\times$ \textbf{1.29}} \\

\bottomrule
\end{tabular}
}
\end{table*}
}}

% ================= Table: Task breakdown (Speedup only) =================
\newcommand{\specbenchTable}{{
\begin{table*}[ht]
\centering
\caption{Goodput (tokens/sec) of \sysname and baselines on Spec-Bench on Qwen-3-8B. The best performance is marked in bold. The speedup is measured by comparing the goodput with the normal autoregressive decoding.}
\label{tab:qwen_specbench}
\resizebox{1\linewidth}{!}{
\begin{tabular}{c| c| cc cc cc cc cc cc cc}
\toprule
\textsc{Model} & \textsc{Method}
& \multicolumn{2}{c}{MT Conversation}
& \multicolumn{2}{c}{Question Answering}
& \multicolumn{2}{c}{Summarization}
& \multicolumn{2}{c}{Translation}
& \multicolumn{2}{c}{Math Reasoning}
& \multicolumn{2}{c}{Retrieval Augmented}
& \multicolumn{2}{c}{All}\\
\cmidrule(lr){3-4}\cmidrule(lr){5-6}\cmidrule(lr){7-8}\cmidrule(lr){9-10}
\cmidrule(lr){11-12}\cmidrule(lr){13-14}\cmidrule(lr){15-16}
& & Tokens/s & Speedup & Tokens/s & Speedup & Tokens/s & Speedup
& Tokens/s & Speedup & Tokens/s & Speedup & Tokens/s & Speedup & Tokens/s & Speedup \\
\midrule

\multirow{5}{*}{\rotatebox[origin=c]{90}{\textbf{Qwen-3-8B}}}
&           Normal
& 939.50 & \spone{$\times$1.00}
& 860.97 & \spone{$\times$1.00}
& 1277.50 & \spone{$\times$1.00}
& 2013.02 & \spone{$\times$1.00}
& 1064.78 & \spone{$\times$1.00}
& 191.60 & \spone{$\times$1.00}
& 697.42 & \spone{$\times$1.00} \\\cmidrule{2-16}
&           EAGLE-3
& 1005.63 & \spup{$\times$1.07}
& 895.25  & \spup{$\times$1.04}
& 1552.62 & \spup{$\times$1.22}
& 2319.01 & \spup{$\times$1.15}
& 1045.28 & \spdown{$\times$0.98}
& 168.21  & \spdown{$\times$0.88}
& 783.25  & \spup{$\times$1.12} \\
&           NGram
& 581.26 & \spdown{$\times$0.62}
& 351.28 & \spdown{$\times$0.41}
& 1082.57 & \spdown{$\times$0.85}
& 2166.37 & \spup{$\times$1.08}
& 959.38 & \spdown{$\times$0.90}
& 263.11 & \spup{$\times$1.37}
& 545.39 & \spdown{$\times$0.78} \\
&           SuffixDecoding
& 1007.08 & \spup{$\times$1.07}
& 873.12  & \spup{$\times$1.01}
& 1213.63 & \spdown{$\times$0.95}
& 2141.79 & \spup{$\times$1.06}
& 1121.48 & \spup{$\times$1.05}
& 212.19 & \spup{$\times$1.10}
& 712.79 & \spup{$\times$1.02} \\\cmidrule{2-16}
&           \sysname
& \textbf{1128.09} & \spup{$\times$ \textbf{1.20}}
& \textbf{929.85} & \spup{$\times$ \textbf{1.08}}
& \textbf{1481.90} & \spup{$\times$ \textbf{1.16}}
& \textbf{2294.12} & \spup{$\times$ \textbf{1.14}}
& \textbf{1075.43} & \spup{$\times$ \textbf{1.01}}
& \textbf{267.14} & \spup{$\times$ \textbf{1.40}}
& \textbf{796.01} & \spup{$\times$ \textbf{1.14}} \\
\bottomrule
\end{tabular}
}
\end{table*}
}}


\newcommand{\sensitivityTable}{{
\begin{table}[t]
\centering
\caption{Goodput of variants of \sysname and baselines on Deep Research Agent and Code Generation Agent workloads. The best performance is marked in bold. The speedup is measured by comparing the goodput with the normal autoregressive decoding.}
\label{tab:sensitivity}
\resizebox{\linewidth}{!}{
\begin{tabular}{c|c|cc|cc}
\toprule
\textsc{Model} & \textsc{Method}
& \multicolumn{2}{c|}{Deep Research Agent}
& \multicolumn{2}{c}{Code Generation Agent} \\
\cmidrule(lr){3-4}\cmidrule(lr){5-6}
& & Tokens/s & Speedup & Tokens/s & Speedup \\
\midrule
\multirow{7}{*}{\rotatebox[origin=c]{90}{\textsc{\textbf{Qwen-3-8B}}}}
& Normal         & 617.41 & \spone{$\times$ 1.00}                 & 636.09 & \spone{$\times$ 1.00} \\ \cmidrule(lr){2-6}
& EAGLE-3         & 605.14 & \spdown{$\times$ 0.98} & 539.01 & \spdown{$\times$ 0.85} \\
& NGram          & 509.62 & \spdown{$\times$ 0.83} & 422.58 & \spdown{$\times$ 0.66} \\
& SuffixDecoding & 598.74 & \spdown{$\times$ 0.97} & 579.21 & \spdown{$\times$ 0.91} \\
\cmidrule(lr){2-6}
& \textsc{\sysname\texttt{(S)}}
& \textbf{647.12}~\spup{($\uparrow 9\%$)}
& \spup{$\times$ \textbf{1.05}}
& \textbf{795.54}~\spup{($\uparrow 37\%$)}
& \spup{$\times$ \textbf{1.25}} \\
& \textsc{\sysname\texttt{(R)}}
& \textbf{623.33}~\spup{($\uparrow 3\%$)}
& \spup{$\times$ \textbf{1.01}}
& \textbf{702.28}~\spup{($\uparrow 21\%$)}
& \spup{$\times$ \textbf{1.10}} \\  \cmidrule{2-6}
& \textsc{\sysname\texttt{(N)}}
& \textbf{661.08}~\spup{($\uparrow 10\%$)}
& \spup{$\times$ \textbf{1.07}}
& \textbf{809.12}~\spup{($\uparrow 40\%$)}
& \spup{$\times$ \textbf{1.27}} \\   \cmidrule{2-6}
& \textsc{\sysname}
& \textbf{661.43}~\spup{($\uparrow 10\%$)}
& \spup{$\times$ \textbf{1.07}}
& \textbf{828.10}~\spup{($\uparrow 43\%$)}
& \spup{$\times$ \textbf{1.30}} \\
\bottomrule
\end{tabular}
}
\end{table}
}}

\newcommand{\rejTable}{{
\begin{table}[t]
\centering
\caption{Average per-request rejection rate \spup{($\downarrow$)}, defined as number of rejected tokens divided by number of draft tokens, for different speculative decoding algorithms on Reflexion agentic workflow.}
\label{tab:rejection_rate}
\resizebox{1\linewidth}{!}{
\begin{tabular}{c|ccc|c}
\toprule
\textsc{Metric} & EAGLE-3 & NGram & SuffixDecoding  & \sysname\\
\midrule
Rejection Rate & 53.3\% & 88.7\% & 65.9\%  & 26.4\% \spup{($\downarrow 50\%$)}\\
\bottomrule
\end{tabular}
}
\end{table}
}}

\newcommand{\breakdownTable}{{
\begin{table}[t]
\centering
\caption{Execution breakdown of \sysname and baselines on Code Gen agentic Workload. The best performance is marked in bold. }
\label{tab:runtime_breakdown}
\resizebox{\linewidth}{!}{%
\begin{tabular}{l|cccc}
\toprule
\textsc{Time (minutes)} & EAGLE-3 & NGram & SuffixDecoding & \sysname \\
\midrule
Draft        & 5.72  & \textbf{0.68}  & 1.41  & 1.74  \\\midrule
Verification & 34.24 & 39.82 & 33.12 & \textbf{26.11}~\spup{($\downarrow 21\%$)} \\\midrule
Overall      & 39.96 & 40.50 & 34.53 & \textbf{27.85}~\spup{($\downarrow 19\%$)} \\
\bottomrule
\end{tabular}%
}
\end{table}
}}

% \newcommand{\acclenTable}{{
% \begin{table}[t]
% \centering
% \caption{Acceptance length of \sysname under different template settings of Qwen-3-8B on CodeGen agent workload.}
% \label{tab:template_acceptance_length}
% \resizebox{0.9\linewidth}{!}{%
% \begin{tabular}{l|ccc}
% \toprule
% \textbf{Setting} & \textbf{Template} & \textbf{Non-template} & \textbf{Overall} \\
% \midrule
% Acceptance Len. & 4.03 & 3.72 & 3.74 \\
% \bottomrule
% \end{tabular}%
% }
% \end{table}
% }}

\newcommand{\adaTable}{{
\centering
\caption{Speedup of \sysname compared to AdaSpec on Qwen-3-8B across agent workloads.}
\label{tab:adaspec_speedup}
\resizebox{\linewidth}{!}{%
\begin{tabular}{l|cccc}
\toprule
\textbf{Qwen-3-8B} & \textbf{DeepResearch} & \textbf{CodeGen} & \textbf{GAIA} & \textbf{SWE-Bench} \\
\midrule
AdaSpec       & 0.91$\times$ & 1.02$\times$ & 0.99$\times$ & 1.03$\times$ \\
AgentSpec (R) & 1.01$\times$ & 1.10$\times$ & 1.05$\times$ & 1.20$\times$ \\
AgentSpec     & \textbf{1.07$\times$} & \textbf{1.30$\times$} & \textbf{1.07$\times$} & \textbf{1.31$\times$} \\
\bottomrule
\end{tabular}%
}
}}

\newcommand{\GPUTable}{{
\centering
\caption{Speedup of \sysname on different GPU platforms on GPT-OSS-20B across agent workloads.}
\label{tab:hardware_speedup}
\resizebox{\linewidth}{!}{%
\begin{tabular}{l|cccc}
\toprule
\textbf{Speedup} & \textbf{CodeGen} & \textbf{DeepResearch} & \textbf{SWE-Bench} & \textbf{GAIA} \\
\midrule
1$\times$A100 80G & 2.02$\times$ & 1.48$\times$ & 1.69$\times$ & 1.19$\times$ \\
1$\times$H100 80G & 2.28$\times$ & 1.67$\times$ & 1.92$\times$ & 1.54$\times$ \\
\bottomrule
\end{tabular}%
}
}}
